% 00-final-sample.tex – TEMPORARY final design sample (Step 13). Placeholder text only; not book content.
\ChapterMission{המשימה: להראות שהבעיה אמיתית, מדידה ומתרחשת כל בוקר}
\ChapterIcon{pedestrian}
\chapter{דוגמת עיצוב סופית: פרק בשלב הבעיה}
\begin{scenebox}
שבע וחצי בבוקר. האור האדום דולק, והכביש ריק לגמרי. אישה אחת עומדת לבדה על המדרכה,
מביטה ימינה ושמאלה, ומעבירה את משקלה מרגל לרגל. אין איש לידה שיראה אותה.
היא יורדת מהמדרכה.
\end{scenebox}

זהו טקסט לדוגמה שנועד לבדוק את העיצוב הסופי בלבד, ואינו חלק מתוכן הדוח.\sidenote{הערת שוליים לדוגמה בשוליים החיצוניים.}
הפסקה משלבת עברית עם מונחים באנגלית, כמו \en{Smart Pedestrian Crossing} ו-\en{countdown signal},
וכן מספרים: \pct{13.5} מההולכים שהגיעו במופע האדום חצו באדום, מתוך \num{563} תצפיות במופע זה
ו-\num{1,392} תצפיות בסך הכול \cite{rosenbloom2009}. רווח הסמך לדוגמה הוא \range{0.51}{0.78}.

\begin{statrow}[3]
  \stattile[SignalRed]{\pct{13.5}}{מהמגיעים במופע האדום חצו באדום}
  \stattile[SignalAmber]{\num{1,392}}{הולכי רגל נצפו בצומת אחד בשעת בוקר}
  \stattile[SignalGreen]{\num{0.63}}{יחס הסיכויים לכל ממתין נוסף (מחושב מחדש)}
\end{statrow}

\begin{pullquote}[ציטוט לדוגמה, לא מתוך המקור]
כשאין איש על המדרכה, הרמזור הוא רק הצעה.
\end{pullquote}

\section{סעיף לדוגמה: תיבות, טבלה, נוסחה ותמונה}
משפט מעבר קצר מציג את רכיבי העמוד. טבלה~\ref{tab:final} מציגה נתונים שעובדו מן המקור, משוואה~\eqref{eq:final}
מציגה את המודל, ואיור~\ref{fig:photo} מציג צילום של המחברים.

\begin{findingbox}
ממצא לדוגמה: ככל שממתינים יותר אנשים על המדרכה, כך יורד הסיכוי לחצייה באדום.
\end{findingbox}
\begin{cautionbox}
סוגיה לדוגמה: עמודת \en{$e^{b}$} בטבלה המקורית אינה תואמת את \en{$\exp(b)$} בשני משתנים.
\end{cautionbox}

\begin{table}[ht]
  \centering
  \caption{טבלה לדוגמה: חצייה באדום בקרב המגיעים במופע האדום, לפי מגדר (עובד מתוך \cite{rosenbloom2009})}
  \label{tab:final}
  \begin{tabular}{lccc}
    \toprule
    \headrow \thead{מגדר} & \thead{הגיעו באדום} & \thead{חצו באדום} & \thead{שיעור} \\
    \midrule
    נשים & \num{330} & \num{28} & \pct{8.5} \\
    גברים & \num{233} & \num{48} & \pct{20.6} \\
    סך הכול & \num{563} & \num{76} & \pct{13.5} \\
    \bottomrule
  \end{tabular}
\end{table}

\begin{equation}
  \ln\frac{p}{1-p} = \beta_0 + \beta_1 x_{\hebmath{מגדר}} + \beta_2 x_{\hebmath{ממתינים}} + \dots + \beta_k x_k
  \label{eq:final}
\end{equation}

\authorphoto[80mm]{figures/photo-crosswalk.jpg}{רמזור הולכי רגל עם תצוגת ספירה לאחור מעל מעבר חציה}{fig:photo}
\crosswalkrule
